\subsection{2.4 Dữ liệu nguồn}

\begin{frame}{Phạm vi: dữ liệu 01/2017–08/2018 chiếm 99,6\% số đơn nên được dùng cho xu hướng}
\muc{2.4 · Dữ liệu nguồn}
\begin{columns}[T]\begin{column}{0.4\textwidth}
\begin{itemize}
  \item 2016 chỉ 329 đơn; 09–10/2018 chỉ 20 đơn
\end{itemize}
\end{column}\begin{column}{0.58\textwidth}
\begin{center}\anh[0.74\textheight]{3-3-1_orders_per_month.png}\end{center}
\end{column}\end{columns}

\note{[S32] Bộ lọc InWindow trong Power BI cố định cửa sổ này cho DB1–DB4.}
\end{frame}

\begin{frame}{Hai hệ thống nguồn có 11 bảng, 121 MB và 1.559.764 dòng}
\muc{2.4 · Dữ liệu nguồn}
\begin{itemize}
  \item oltp 9 bảng 1.550.922 dòng; crm 2 bảng 8.842 dòng
\end{itemize}
\begin{center}\adjustbox{max width=\linewidth,max height=0.5\textheight}{\scriptsize\begin{tabular}{llrrr}
\toprule
\textbf{Bảng nguồn} & \textbf{File CSV} & \textbf{Số dòng} & \textbf{Số cột} & \textbf{MB} \\
\midrule
oltp.product\_category\_name\_translation & product\_category\_name\_translation.csv & 71 & 2 & 0,00 \\
oltp.products & olist\_products\_dataset.csv & 32.951 & 9 & 2,27 \\
oltp.sellers & olist\_sellers\_dataset.csv & 3.095 & 4 & 0,17 \\
oltp.geolocation & olist\_geolocation\_dataset.csv & 1.000.163 & 6 & 58,44 \\
oltp.customers & olist\_customers\_dataset.csv & 99.441 & 5 & 8,62 \\
oltp.orders & olist\_orders\_dataset.csv & 99.441 & 8 & 16,84 \\
oltp.order\_items & olist\_order\_items\_dataset.csv & 112.650 & 7 & 14,72 \\
oltp.order\_payments & olist\_order\_payments\_dataset.csv & 103.886 & 5 & 5,51 \\
oltp.order\_reviews & olist\_order\_reviews\_dataset.csv & 99.224 & 7 & 13,78 \\
crm.mql & olist\_marketing\_qualified\_leads\_dataset.csv & 8.000 & 4 & 0,67 \\
crm.closed\_deals & olist\_closed\_deals\_dataset.csv & 842 & 14 & 0,16 \\
Tổng (11 bảng) &  & 1.559.764 &  & 121,18 \\
\bottomrule
\end{tabular}}\end{center}

\note{[S33] Vượt tiêu chí kích thước dữ liệu của môn (trên 100 MB).}
\end{frame}

\begin{frame}{Bảng toạ độ chiếm 64\% số dòng nhưng chỉ đại diện 19.015 mã bưu chính}
\muc{2.4 · Dữ liệu nguồn}
\begin{itemize}
  \item 1.000.163 dòng geolocation
\end{itemize}
\khungchen[0.3\textheight]{CHÈN HÌNH / BẢNG: bảng nguồn (làm nổi dòng geolocation)}

\note{[S34] Đây là bảng cần xử lý nhiều nhất ở ODS.}
\end{frame}

\begin{frame}{Luồng dữ liệu và tần suất cập nhật: nạp ban đầu 45.430 đơn rồi 11 đợt hằng tháng}
\muc{2.4 · Dữ liệu nguồn}
\begin{columns}[T]\begin{column}{0.4\textwidth}
\begin{itemize}
  \item Tổng 286.144 bản ghi mới qua 11 đợt
\end{itemize}
\end{column}\begin{column}{0.58\textwidth}
\begin{center}\anh[0.74\textheight]{3-4-1_pipeline.png}\end{center}
\end{column}\end{columns}

\note{[S37] Dữ liệu Kaggle là ảnh tĩnh nên nhóm mô phỏng cập nhật.}
\end{frame}

\begin{frame}{Trung bình 163 đơn mỗi ngày, đỉnh 1.176 đơn ngày Black Friday: lô hằng ngày là đủ}
\muc{2.4 · Dữ liệu nguồn}
\begin{columns}[T]\begin{column}{0.4\textwidth}
\begin{itemize}
  \item Cửa sổ 01/2017--08/2018: 163 đơn, 185 dòng hàng, 13 MQL, 1,3 deal mỗi ngày
  \item Bảng tần suất đề xuất công nghệ: lô hằng ngày bằng procedure PostgreSQL; gần thời gian thực thì CDC + Kafka (chương 5)
\end{itemize}
\end{column}\begin{column}{0.58\textwidth}
\begin{center}\anh[0.74\textheight]{2-4_tan_suat_cap_nhat.png}\end{center}
\end{column}\end{columns}

\note{[S127] Biểu đồ tần suất cập nhật theo thời gian và bảng công nghệ xử lý đề xuất.}
\end{frame}

\begin{frame}{Cây phân loại dữ liệu chia 11 bảng thành 4 nhóm: chủ, giao dịch, tương tác, tham chiếu}
\muc{2.4 · Dữ liệu nguồn}
\begin{columns}[T]\begin{column}{0.4\textwidth}
\begin{itemize}
  \item Mỗi bảng kèm hai cột tiêu biểu và giá trị thật
  \item Phân loại thêm theo kiểu dữ liệu, độ nhạy cảm (định danh, vị trí) và tính biến động
\end{itemize}
\end{column}\begin{column}{0.58\textwidth}
\begin{center}\anh[0.74\textheight]{2-4_data_taxonomy.png}\end{center}
\end{column}\end{columns}

\note{[S128] Data taxonomy của 11 bảng nguồn.}
\end{frame}

\begin{frame}{Data flow hiện trạng: mũi tên ghi thông tin được truyền, hai hệ thống nguồn tách rời}
\muc{2.4 · Dữ liệu nguồn}
\begin{columns}[T]\begin{column}{0.4\textwidth}
\begin{itemize}
  \item Thông tin đơn, trạng thái giao, thanh toán, điểm đánh giá, lead
  \item Marketing funnel (SDR, SR) tách khỏi hệ thống bán hàng
\end{itemize}
\end{column}\begin{column}{0.58\textwidth}
\begin{center}\anh[0.74\textheight]{2-4_data_flow_hien_trang.png}\end{center}
\end{column}\end{columns}

\note{[S129] Hiện trạng: luồng thông tin giữa các hệ thống CNTT (bản nháp, nhóm vẽ lại).}
\end{frame}

\begin{frame}{Data flow đề xuất: nguồn, staging, ODS, kho, mart và BI, có nhật ký ETL giám sát}
\muc{2.4 · Dữ liệu nguồn}
\begin{columns}[T]\begin{column}{0.4\textwidth}
\begin{itemize}
  \item Bảng lớp, thành phần, công nghệ, tần suất, ghi chú ở báo cáo
  \item Điều phối bằng một lệnh Python, biến đổi bằng PL/pgSQL
\end{itemize}
\end{column}\begin{column}{0.58\textwidth}
\begin{center}\anh[0.74\textheight]{2-4_data_flow_de_xuat.png}\end{center}
\end{column}\end{columns}

\note{[S130] Đề xuất: luồng nguồn đến BI.}
\end{frame}

\subsection{2.5 ERD OLTP}

\begin{frame}{ERD hệ thống bán hàng: chín bảng nối quanh đơn hàng}
\muc{2.5 · ERD OLTP}
\begin{columns}[T]\begin{column}{0.4\textwidth}
\begin{itemize}
  \item orders $\rightarrow$ items, payments, reviews; customers, sellers, products
\end{itemize}
\end{column}\begin{column}{0.58\textwidth}
\begin{center}\anh[0.74\textheight]{3-2-4_erd_oltp.png}\end{center}
\end{column}\end{columns}

\note{[S35] \texttt{customer\_\allowbreak{}id} được sinh riêng cho từng đơn nên cần \texttt{customer\_\allowbreak{}unique\_\allowbreak{}id} để nhận ra khách thật.}
\end{frame}

\begin{frame}{Các cột được gạch chân trên ERD là cột đi sang kho, kèm ghi chú mapping}
\muc{2.5 · ERD OLTP}
\begin{itemize}
  \item Ví dụ: \texttt{order\_\allowbreak{}purchase\_\allowbreak{}timestamp} $\rightarrow$ \texttt{purchase\_\allowbreak{}date\_\allowbreak{}key}, \texttt{price} $\rightarrow$ \texttt{price}
\end{itemize}
\khungchen[0.3\textheight]{CHÈN HÌNH / BẢNG: ERD đã gạch chân (vẽ lại)}

\note{[S36] Bảng mapping đầy đủ 108 dòng ODS $\rightarrow$ kho nằm ở phụ lục.}
\end{frame}

\begin{frame}{ERD tô màu theo nhóm dữ liệu: 65 trong 71 cột nguồn được gạch chân vì đi sang OLAP}
\muc{2.5 · ERD OLTP}
\begin{columns}[T]\begin{column}{0.4\textwidth}
\begin{itemize}
  \item Gạch chân suy ra từ bảng mapping nguồn $\rightarrow$ ODS $\rightarrow$ kho
  \item 6 cột không dùng: geolocation.id, 3 thuộc tính sản phẩm, 2 thuộc tính khai báo của deal
\end{itemize}
\end{column}\begin{column}{0.58\textwidth}
\begin{center}\anh[0.74\textheight]{2-5_erd_oltp_gach_chan.png}\end{center}
\end{column}\end{columns}

\note{[S131] ERD nháp tô màu và gạch chân cột dùng sang OLAP.}
\end{frame}

\begin{frame}{Luồng OLTP sang OLAP cho thấy bảng nguồn nào nuôi dimension nào và fact nào}
\muc{2.5 · ERD OLTP}
\begin{columns}[T]\begin{column}{0.4\textwidth}
\begin{itemize}
  \item Bảng tổng hợp mapping theo bảng nguồn ở báo cáo mục 2.5
  \item Mapping đầy đủ từng cột ở Phụ lục C
\end{itemize}
\end{column}\begin{column}{0.58\textwidth}
\begin{center}\anh[0.74\textheight]{2-5_data_flow_oltp_olap.png}\end{center}
\end{column}\end{columns}

\note{[S132] Sơ đồ luồng dữ liệu OLTP sang OLAP.}
\end{frame}

